\subsection{Orbital mechanics and the top boundary --- full picture with diagrams}
\label{sec:orbit-top-bc}

The \textbf{top boundary} is where the Moon meets space. Unlike Earth, there is
no atmosphere to scatter or absorb sunlight, and no wind to carry heat away.
Everything that heats or cools the surface comes from:
\begin{enumerate}[leftmargin=1.4em,itemsep=2pt]
  \item \textbf{Direct sunlight} (orbital geometry sets how strong it is at each moment),
  \item \textbf{Thermal radiation} the surface emits to the cold sky,
  \item \textbf{Conduction} into the regolith below.
\end{enumerate}
This subsection explains the \textbf{orbital part} (where does $F_\odot(t)$ come
from?) and ties it to the \textbf{surface energy balance} that closes the top BC.

\subsubsection{The Moon in orbit: what ``one day'' means}

The Moon rotates once per orbit around Earth, so the same face always points
toward Earth. From a point on the lunar surface, ``day'' and ``night'' are set by
the Sun, not by watching Earth.

\begin{itemize}[leftmargin=1.4em,itemsep=2pt]
  \item \textbf{Synodic month} $P \approx 29.53$~days: time from one local noon
    (Sun highest) to the next. This is the period used in the code
    (\code{T\_LUNAR} $\approx 2.55\times10^6$~\si{s}).
  \item \textbf{Subsolar point}: the spot on the Moon where the Sun is directly
    overhead ($\theta_\odot = 0$). It marches around the equator once per synodic month.
  \item \textbf{Site latitude} $\phi$: Apollo~15 is at $26.1^\circ$N, Apollo~17 at
    $20.2^\circ$N. The Sun never passes directly overhead there; at local ``noon'' the
    Sun is lower in the sky than at the equator.
\end{itemize}

\begin{plainenglish}
Imagine standing on the Moon at an Apollo site. The Sun rises, tracks across the
sky at a height that depends on your latitude, sets, and you wait through a long
night (${\sim}$14 Earth days) until it rises again. The model does not simulate
every crater shadow or mountain --- it uses a smooth ``average day'' for that latitude.
\end{plainenglish}

% ---- Figure: Sun--Moon--surface geometry --------------------------------
\begin{figure}[htbp]
\centering
\begin{tikzpicture}[
  x=1cm, y=1cm,
  font=\small,
  sun/.style={circle, fill=yellow!80!orange, draw=orange!70!black, line width=0.6pt, minimum size=1.1cm},
  moon/.style={circle, fill=softgrey!40, draw=dim, line width=0.8pt},
  ray/.style={-{Stealth[length=2.5mm]}, line width=0.7pt, color=orange!85!black},
  dashedline/.style={dashed, color=dim, line width=0.5pt},
  label/.style={font=\footnotesize, text=char}
]

% Sun
\node[sun] (sun) at (-5.5, 2.2) {};
\node[label, below=2pt of sun] {Sun};

% Moon (cross-section)
\draw[moon] (0,0) circle (2.2);
\fill[softgrey!25] (0,0) circle (2.2);

% Local horizontal at site (latitude phi)
\coordinate (C) at (0,0);
\coordinate (S) at ({2.2*cos(55)},{2.2*sin(55)});  % site on surface, phi ~ 35 deg from equator visually
\draw[line width=1.2pt, color=char] (C) -- (S);
\node[label, anchor=west] at ($(S)+(0.15,0.05)$) {regolith surface};

% Vertical at site (local zenith)
\draw[dashedline] (S) -- ++(0, 1.8) node[above, label] {local vertical (zenith)};

% Sun ray to site
\draw[ray] (sun.east) .. controls (-2.5, 2.0) and (-0.8, 1.6) .. (S);
\node[label, fill=paper, inner sep=1pt] at (-2.2, 2.55) {sunlight};

% Solar zenith angle theta
\draw[{Stealth[length=2mm]}-, line width=0.6pt, color=teal] (0.35, 0) arc (0:55:0.55);
\node[label, text=teal] at (0.95, 0.22) {$\theta_\odot$};

% Latitude phi from equator
\draw[dashedline] (-2.2, 0) -- (2.2, 0);
\draw[{Stealth[length=2mm]}-, line width=0.6pt, color=forest] (0.7, 0) arc (0:55:0.7);
\node[label, text=forest] at (1.35, -0.15) {$\phi$};

% Equator hint
\node[label, text=dim] at (-2.5, -0.35) {lunar equator (plane of subsolar march)};

% Orbit hint
\draw[-{Stealth[length=2mm]}, line width=0.5pt, color=plum, dashed]
  (2.5, 0.3) arc (10:350:2.5 and 0.35);
\node[label, text=plum, anchor=west] at (2.9, 0.3) {Moon orbit (schematic)};

% Annotation box
\node[draw=teal!60, fill=teal!5, rounded corners=2pt, align=left, font=\footnotesize,
      text width=4.2cm, anchor=north west] at (-5.8, -2.6) {
  \textbf{Zenith angle} $\theta_\odot$: angle between\\
  local vertical and the Sun.\\
  Smaller $\theta_\odot$ $\Rightarrow$ stronger heating.
};

\end{tikzpicture}
\caption{\textbf{Sun--Moon--surface geometry.} At a site at latitude $\phi$, the
  Sun is not overhead at ``noon'' unless $\phi=0$. The \textbf{solar zenith angle}
  $\theta_\odot$ sets how much solar power hits each square metre of horizontal
  ground (cosine law). The retrieval uses an idealised cycle
  $\cos\theta_\odot(t)=\cos\phi\cos(2\pi t/P)$; full SPICE ephemeris lives in
  \file{lunar/ephem.py} but is not required for the \Kdstar{} fit.}
\label{fig:sun-moon-geometry}
\end{figure}

\subsubsection{From geometry to insolation: the cosine law}

The solar constant $S_0 \approx \SI{1361}{W/m^2}$ is the flux in a beam
perpendicular to the Sun. On a \textbf{horizontal} patch of ground, only the
component normal to the beam counts:
\begin{equation}
  F_\odot(t) = \max\!\bigl(0,\; S_0 \cos\theta_\odot(t)\bigr).
  \label{eq:insolation-teach}
\end{equation}
When the Sun is low ($\theta_\odot$ large), $\cos\theta_\odot$ is small and
heating is weak. At night, $\cos\theta_\odot < 0$ and the $\max(0,\cdot)$ clips
insolation to zero.

\paragraph{Idealised synodic model used in production.}
For each site latitude $\phi$,
\begin{equation}
  \cos\theta_\odot(t) = \cos\phi \,\cos\!\left(\frac{2\pi t}{P}\right),
  \label{eq:cos-zenith-ideal}
\end{equation}
with $P = \code{T\_LUNAR}$. This is the same formula as in the manuscript
(Eq.~\eqref{eq:insolation-teach}). It assumes:
\begin{itemize}[leftmargin=1.4em,itemsep=1pt]
  \item smooth variation over one lunation (no libration jitter at the hour scale),
  \item no solar \textbf{declination} tilt ($\pm 1.5^\circ$ neglected),
  \item no orbital \textbf{eccentricity} changing $S_0$ ($<0.1\%$ effect on cycle mean).
\end{itemize}
Those corrections matter for \emph{phase-resolved} work; this project compares
\textbf{time-mean} temperatures over a stability window, so the simplified cycle is
adequate (see manuscript \S``Surface solar geometry'').


% ---- Figure: insolation over one lunation --------------------------------
\begin{figure}[htbp]
\centering
\begin{tikzpicture}[x=0.42cm, y=0.00035cm, font=\small]
  % Axes
  \draw[-{Stealth}, line width=0.6pt] (0,0) -- (52,0) node[right, font=\footnotesize] {$t$ (days)};
  \draw[-{Stealth}, line width=0.6pt] (0,0) -- (0,4200)
    node[above, font=\footnotesize] {$F_\odot$ (\si{W/m^2})};
  \draw[color=dim!60, line width=0.3pt] (0,0) -- (0,4200);
  \draw[color=dim!60, line width=0.3pt] (7.38,0) -- (7.38,4200);
  \draw[color=dim!60, line width=0.3pt] (14.76,0) -- (14.76,4200);
  \draw[color=dim!60, line width=0.3pt] (22.14,0) -- (22.14,4200);
  \draw[color=dim!60, line width=0.3pt] (29.53,0) -- (29.53,4200);
  \node[below, font=\tiny, text=dim] at (0,0) {0};
  \node[below, font=\tiny, text=dim] at (7.38,0) {7.4};
  \node[below, font=\tiny, text=dim] at (14.76,0) {14.8};
  \node[below, font=\tiny, text=dim] at (22.14,0) {22.1};
  \node[below, font=\tiny, text=dim] at (29.53,0) {29.5};
  \node[below, font=\footnotesize, text=dim] at (26, -28)
    {one synodic month $P \approx 29.5$~d};

  % cos curve for phi = 26 deg (peak ~ 1223 W/m2); day half only
  \draw[line width=1.1pt, color=coral, domain=0:14.76, samples=80]
    plot (\x, {1361*0.899*cos(360*\x/29.53)});

  % Night shading
  \fill[dim!12] (14.76,0) rectangle (29.53,4200);

  \node[font=\footnotesize, text=dim, rotate=90] at (-3.5, 600) {night ($F_\odot=0$)};
  \node[font=\footnotesize, text=coral] at (7.5, 1150) {daytime insolation};
  \node[font=\footnotesize, align=center] at (38, 3500) {Apollo~15-scale\\$\phi=26.1^\circ$N\\peak $\lesssim S_0\cos\phi$};

  % S0 reference
  \draw[dashed, color=teal, line width=0.6pt] (0,1361) -- (50,1361);
  \node[font=\tiny, text=teal, anchor=west] at (50.5, 1361) {$S_0$};

\end{tikzpicture}
\caption{\textbf{Idealised insolation} $F_\odot(t)$ for northern mid-latitude
  ($\cos\phi \approx 0.90$). Sunlit half of the lunation: flux follows a cosine;
  night: clipped to zero. Peak is below $S_0$ because the Sun never reaches the
  zenith at $\phi > 0$. Code: \code{insol = S0 * cos\_lat * np.maximum(0, np.cos(phase))}
  in \file{retrieve\_kd.py}.}
\label{fig:insolation-lunation}
\end{figure}

\subsubsection{Top boundary: surface energy balance (control volume)}

Take a thin control volume \emph{at} the surface --- the top face sees space, the
bottom face touches the first regolith cell. In steady balance at each instant,
Eq.~\eqref{eq:bc-top-teach} applies (derived in \autoref{sec:bc-full}), with
insolation written explicitly as $F_\odot(t)$ instead of $S(t)$:
\[
  (1-A)\,F_\odot(t) \;=\; \varepsilon\sigma T_s(t)^4 \;+\; q_{\mathrm{cond}}(t).
\]
Here $q_{\mathrm{cond}} = -K\,\partial T/\partial z\big|_{z=0}$ is the downward
conductive flux (positive when heat enters the ground).

% ---- Figure: surface energy balance --------------------------------------
\begin{figure}[htbp]
\centering
\begin{tikzpicture}[
  x=1cm, y=1cm, font=\small,
  arr/.style={-{Stealth[length=2.8mm]}, line width=0.85pt},
  thinarr/.style={-{Stealth[length=2mm]}, line width=0.6pt}
]

% Regolith column
\fill[softgrey!35] (-1.2, -2.8) rectangle (1.2, 0);
\draw[line width=0.8pt, color=char] (-1.2, -2.8) rectangle (1.2, 0);
\foreach \y in {-0.5,-1.0,-1.5,-2.0,-2.5} {
  \draw[dim!40, line width=0.3pt] (-1.2,\y) -- (1.2,\y);
}
\node[font=\footnotesize, text=dim] at (0, -1.4) {regolith ($z>0$)};

% Surface line
\draw[line width=1.4pt, color=char] (-1.5, 0) -- (1.5, 0);
\node[anchor=west, font=\footnotesize] at (1.55, 0) {$z=0$ (surface)};

% Sun / insolation from upper left
\draw[arr, color=orange!85!black] (-2.8, 2.2) -- (-0.3, 0.35);
\node[align=left, font=\footnotesize, fill=paper, inner sep=2pt] at (-3.5, 2.5) {
  incident solar\\$F_\odot(t)$
};
\draw[thinarr, color=orange!60!black] (-1.0, 1.5) -- (-0.2, 0.25);
\node[align=left, font=\footnotesize, text=orange!70!black] at (-2.2, 1.0) {
  absorbed\\$(1-A)F_\odot$
};

% Reflected (small arrow back)
\draw[thinarr, color=dim, dashed] (-0.1, 0.2) -- (-1.2, 1.0);
\node[font=\tiny, text=dim] at (-1.5, 1.15) {reflected $AF_\odot$};

% Outgoing IR upward
\draw[arr, color=coral] (0.2, 0.15) -- (0.2, 2.0);
\node[align=left, font=\footnotesize, text=coral, anchor=west] at (0.45, 2.0) {
  emitted thermal radiation\\$\varepsilon\sigma T_s^4$
};

% Conduction downward
\draw[arr, color=teal] (0, -0.05) -- (0, -1.2);
\node[align=center, font=\footnotesize, text=teal, anchor=west] at (0.25, -0.7) {
  conduction\\$q_{\mathrm{cond}}=-K\partial T/\partial z$
};

% Sky
\fill[teal!4] (-3.2, 0.05) rectangle (3.2, 2.8);
\node[font=\footnotesize, text=dim] at (0, 2.45) {space / sky ($T_{\mathrm{space}}\approx 0$)};

% Ts label
\node[circle, fill=coral!15, draw=coral, inner sep=1pt, font=\footnotesize\bfseries] at (0, 0) {$T_s$};

% Balance equation box
\node[draw=char!50, fill=paper, rounded corners=3pt, align=center, font=\footnotesize,
      text width=5.8cm] at (0, -3.55) {
  Top BC (each timestep): \quad
  $(1-A)F_\odot = \varepsilon\sigma T_s^4 + q_{\mathrm{cond}}$\\[3pt]
  \code{surface\_energy\_balance\_residual(Ts)} = 0
  \quad$\Rightarrow$\quad Newton finds $T_s$
};

\end{tikzpicture}
\caption{\textbf{Top boundary control volume.} Absorbed sunlight (after albedo $A$)
  must equal the sum of infrared emission and heat conducted into the regolith.
  $T_s$ is not stored as a fixed number --- it is \textbf{solved} each hour so the
  balance holds. This is what \code{\_solve\_surface\_newton} does in
  \file{solver.py}.}
\label{fig:surface-energy-balance}
\end{figure}

\begin{analogy}
The surface is like a pan on a stove that also glows in the dark. Sunlight pours
in (when it's day), the pan radiates heat to the cold kitchen (space), and the
handle (regolith) conducts heat downward. At every instant the pan's top temperature
adjusts until those three flows balance.
\end{analogy}

\subsubsection{How orbital forcing connects to the heat equation}

The chain from orbit to interior is:
\begin{center}
\begin{tikzpicture}[
  node distance=0.35cm and 0.5cm,
  box/.style={draw, rounded corners=2pt, fill=paper, align=center,
              font=\footnotesize, inner sep=5pt, minimum width=2.6cm},
  arr/.style={-{Stealth[length=2mm]}, line width=0.55pt, color=char}
]
\node[box, fill=orange!8] (orbit) {Orbit + latitude $\phi$\\$\Rightarrow\cos\theta_\odot(t)$};
\node[box, fill=coral!8, right=of orbit] (insol) {$F_\odot(t)=S_0\max(0,\cos\theta_\odot)$};
\node[box, fill=teal!8, right=of insol] (bc) {Top BC:\\$(1-A)F_\odot=\varepsilon\sigma T_s^4+q_{\mathrm{cond}}$};
\node[box, fill=forest!8, right=of bc] (pde) {Interior PDE:\\$\rho c_p\partial T/\partial t=\partial(K\partial T/\partial z)/\partial z$};
\node[box, fill=plum!8, below=of bc] (newton) {Newton $\to T_s$ each step};
\draw[arr] (orbit) -- (insol);
\draw[arr] (insol) -- (bc);
\draw[arr] (bc) -- (pde);
\draw[arr] (bc) -- (newton);
\draw[arr] (newton) |- (pde);
\end{tikzpicture}
\end{center}

Insolation is \textbf{prescribed} (known input from orbit model). Skin temperature
$T_s$ is \textbf{unknown} until the top BC is satisfied. Once $T_s$ is known, it
feeds the first interior cell and Crank--Nicolson advances $T(z,t)$ for all deeper
layers.

\subsubsection{Bottom boundary (for completeness)}

\begin{figure}[htbp]
\centering
\begin{tikzpicture}[x=1cm, y=1cm, font=\small,
  arr/.style={-{Stealth[length=2.5mm]}, line width=0.75pt}]
\fill[softgrey!35] (-1.2, 0) rectangle (1.2, 2.5);
\draw[line width=0.8pt] (-1.2, 0) rectangle (1.2, 2.5);
\node[font=\footnotesize, text=dim] at (0, 1.2) {regolith};
\draw[line width=1.2pt] (-1.5, 0) -- (1.5, 0);
\node[anchor=west, font=\footnotesize] at (1.55, 0) {$z=z_{\mathrm{max}}$ (5~m)};
\draw[arr, color=coral] (0, -0.1) -- (0, -1.1);
\node[align=center, font=\footnotesize, text=coral, anchor=west] at (0.2, -0.65) {
  basal flux $Q_b$\\from lunar interior
};
\node[draw=char!40, fill=paper, rounded corners=2pt, font=\footnotesize, align=center,
      text width=5.5cm] at (0, -2.0) {
  Bottom BC: \quad $-K\dfrac{\partial T}{\partial z}\Big|_{z_{\mathrm{max}}} = Q_b$\\[2pt]
  Weak interior heat; ${\sim}0.02$~\si{W/m^2} at Apollo sites
};
\end{tikzpicture}
\caption{\textbf{Bottom boundary.} A small steady flux from the Moon's interior.
  No orbital variation --- only the top sees the Sun.}
\label{fig:bottom-bc}
\end{figure}

\begin{headsup}
\textbf{What is \emph{not} in the top BC.} No atmosphere, no wind, no latent heat
(evaporation), no ground heat storage in a separate surface layer. The model is
deliberately minimal: radiative balance + 1-D conduction. That matches airless-body
practice for metre-scale regolith (Hayne et al.\ 2017; Vasavada et al.\ 2012).
\end{headsup}

\subsubsection{Code map: from \code{insolation} to \code{\_step}}

\begin{table}[htbp]
\centering
\small
\begin{tabular}{@{}llp{6.2cm}@{}}
\toprule
Step & Function / variable & Role \\
\midrule
1 & \code{cos\_lat = np.cos(np.radians(lat))} & latitude factor \\
2 & \code{phase = 2*np.pi*t/T\_LUNAR} & lunation phase \\
3 & \code{insol = S0 * cos\_lat * max(0, cos(phase))} & build $F_\odot(t)$ \\
4 & \code{run\_with(..., insolation=insol)} & pass into solver \\
5 & \code{\_solve\_surface\_newton(T\_prev, insol\_k, ...)} & top BC $\to T_s$ \\
6 & \code{\_step(...)} & CN step for interior + bottom $Q_b$ \\
\bottomrule
\end{tabular}
\caption{End-to-end path for orbital forcing and boundary conditions in the retrieval pipeline.}
\label{tab:orbit-bc-code}
\end{table}

For full SPICE-based zenith angles (sub-solar latitude, libration, exact timestamps),
see \file{lunar/ephem.py}. The \Kdstar{} retrieval does not require that path because
it fits \emph{mean} temperatures, not diurnal phase.
